\chapter{Dependencies}
\label{ch:dependencies}

Chapter~\ref{ch:environment} explained \emph{what} each technology does. This chapter is about the dependency declarations themselves: where they live, how versions are constrained and locked, what is installed where, and where the declarations disagree with reality.

\section{Where dependencies are declared}

\begin{center}
\begin{tabularx}{\textwidth}{@{}>{\raggedright\arraybackslash}p{4.2cm}L@{}}
\toprule
\textbf{File} & \textbf{What it controls}\\
\midrule
\code{pyproject.toml}, \code{[project] dependencies} & 14 runtime packages: everything needed to serve a query or build the index.\\
\code{pyproject.toml}, \code{[dependency-groups]} & \code{dev} (9 packages: tests and code quality) and \code{eval} (5 packages: RAGAS and its workarounds).\\
\code{uv.lock} & The complete resolved graph with exact versions and hashes. CI and Docker install with \code{--frozen}, so they use the lockfile exactly as committed. \code{.gitattributes} marks it \code{linguist-generated -diff} so it does not clutter reviews.\\
\code{app/requirements.txt} & What the Hugging Face Space installs. Separate from \code{pyproject.toml}, with major-version ranges and no lockfile.\\
\code{.pre-commit-config.yaml} & Each hook runs in its own environment with its own pinned tool version.\\
\code{Dockerfile} & Installs the runtime set from the lockfile (\code{--no-dev}); the evaluation group is never installed.\\
\bottomrule
\end{tabularx}
\end{center}

\begin{background}[lockfiles and dependency groups]
A version range such as \code{>=1.12,<2.0} says which versions are acceptable; a lockfile records which exact versions were chosen, so every installation from it is identical. Dependency groups (standardised in PEP~735) are named sets of packages that are not part of the published package's requirements, typically used for development and test tools.
\end{background}

\section{Runtime dependencies}

The versions in the \emph{Installed} column are those found in the project's virtual environment during this analysis \tagM.

\begin{longtable}{@{}>{\raggedright\arraybackslash}p{3.4cm}>{\raggedright\arraybackslash}p{2.2cm}>{\raggedright\arraybackslash}p{2.0cm}>{\raggedright\arraybackslash}p{7.7cm}@{}}
\toprule
\textbf{Package} & \textbf{Declared} & \textbf{Installed} & \textbf{Used for, and where}\\
\midrule
\endhead
\code{qdrant-client[fastembed]} & \code{>=1.12,<2.0} & 1.18.0 (fastembed 0.8.0) & Vector store client (\code{store/}, \code{retrieve/}, API health) and local models (\code{llm/provider.py}, \code{retrieve/rerank.py}). Needed everywhere.\\
\code{google-genai} & \code{>=1.0,<2.0} & 1.75.0 & Gemini client, imported lazily inside \code{get_gemini_client}. Needed to answer, grade and decompose; not needed to ingest.\\
\code{pymupdf} & \code{>=1.24,<2.0} & 1.28.0 & PDF parsing and page counts (\code{ingest/parse.py}, \code{ingest/download.py}). Ingestion only.\\
\code{httpx} & \code{>=0.27,<1.0} & 0.28.1 & PDF download (\code{ingest/download.py}). Ingestion only.\\
\code{fastapi} & \code{>=0.115,<1.0} & 0.141.1 & HTTP API (\code{api/}).\\
\code{uvicorn[standard]} & \code{>=0.32,<1.0} & 0.52.0 & Serving the API; used from the command line, not imported.\\
\code{gradio} & \code{>=5.0,<6.0} & 5.50.0 & The demo (\code{app/app.py} only).\\
\code{pydantic} & \code{>=2.9,<3.0} & 2.12.3 & Models throughout the package.\\
\code{pydantic-settings} & \code{>=2.6,<3.0} & 2.14.2 & \code{Settings} (\code{settings.py}).\\
\code{python-dotenv} & \code{>=1.0,<2.0} & 1.2.2 & \code{load_dotenv} in \code{scripts/deploy_space.py}.\\
\code{pyyaml} & \code{>=6.0,<7.0} & 6.0.3 & Profile loading (\code{settings.py}).\\
\code{structlog} & \code{>=24.4,<26.0} & 25.5.0 & Logging configuration (\code{observe/logging.py}).\\
\code{typer} & \code{>=0.15,<1.0} & 0.27.0 & CLI (\code{cli.py}).\\
\code{rich} & \code{>=13.9,<15.0} & 14.3.4 & CLI tables and progress bars.\\
\bottomrule
\end{longtable}

Transitive packages worth knowing by name, because they do real work: \code{onnxruntime} 1.28.0 (model inference), \code{tokenizers} 0.23.1 (the tokenizer the chunker counts with), \code{numpy} 2.4.6, \code{huggingface-hub} 1.25.1 (model downloads, and the deploy script), and \code{starlette} 0.52.1 (under FastAPI) \tagM.

\section{Development and evaluation groups}

\begin{center}
\begin{tabularx}{\textwidth}{@{}>{\raggedright\arraybackslash}p{3.0cm}>{\raggedright\arraybackslash}p{2.6cm}>{\raggedright\arraybackslash}p{2.6cm}L@{}}
\toprule
\textbf{Package} & \textbf{Declared} & \textbf{Installed} & \textbf{Purpose}\\
\midrule
\multicolumn{4}{@{}l}{\textit{Group \code{dev}}}\\
\code{pytest} & \code{>=8.3,<9.0} & 8.4.2 & Test runner\\
\code{pytest-asyncio} & \code{>=0.24,<2.0} & 1.4.0 & Asynchronous tests (\code{asyncio_mode = "auto"})\\
\code{pytest-cov} & \code{>=6.0,<8.0} & 7.1.0 & Coverage, floor 60\%\\
\code{ruff} & \code{>=0.8,<1.0} & 0.16.0 & Linter\\
\code{black} & \code{>=24.10,<26.0} & 25.12.0 & Formatter\\
\code{mypy} & \code{>=1.13,<2.0} & 1.20.2 & Strict type checker\\
\code{pre-commit} & \code{>=4.0,<5.0} & 4.6.1 & Git hooks\\
\code{types-pyyaml} & \code{>=6.0,<7.0} & \code{6.0.12.20260724} & Type stubs for PyYAML\\
\code{respx} & \code{>=0.22,<1.0} & 0.23.1 & HTTP mocking\\
\midrule
\multicolumn{4}{@{}l}{\textit{Group \code{eval}}}\\
\code{ragas} & \code{>=0.4,<0.5} & 0.4.3 & RAG metrics\\
\code{langchain-community} & \code{<0.4} & 0.3.31 & Upstream bug pin (below)\\
\code{openai} & \code{>=1.60,<3.0} & 1.109.1 & Client for Google's OpenAI-compatible endpoint\\
\code{pandas} & \code{>=2.2,<3.0} & 2.3.3 & Not imported by project code\\
\code{datasets} & \code{>=3.0,<5.0} & 4.8.5 & Not imported by project code; required by RAGAS itself\\
\bottomrule
\end{tabularx}
\end{center}

RAGAS 0.4.3 in turn requires, among others, \code{instructor}, \code{langchain}, \code{langchain-core}, \code{langchain-openai}, \code{tiktoken}, \code{networkx} and \code{scikit-network} \tagM. None of these reaches the runtime image, because the evaluation group is installed only where evaluation runs.

\section{Why the unusual constraints exist}

\paragraph{The evaluation group is separate.} \enquote{Evaluation tooling (ragas -> langchain) is deliberately NOT here: it would add \textasciitilde200MB to the deployed image for code that never runs in production} \ev{pyproject.toml:15-17}. The 200\,MB figure is stated, not measured in the repository.

\paragraph{\code{langchain-community<0.4} is an upstream bug pin.} RAGAS 0.4.x imports \code{langchain_community.chat_models.vertexai}, which \code{langchain-community} removed in 0.4.0, so \code{import ragas} fails outright. The comment cites RAGAS issues \#2741 and \#2745 and records that downgrading RAGAS to 0.3.9 fails the same way, because the break is on the LangChain side \ev{pyproject.toml:74-83}.

\paragraph{\code{openai} is a client, not a provider.} RAGAS recognises an asynchronous LLM client by the attribute path \code{chat.completions.create}. Pointing an \code{AsyncOpenAI} client at \code{https://generativelanguage.googleapis.com/v1beta/openai/} satisfies that check while every request still goes to Gemini \ev{pyproject.toml:84-94}.

\paragraph{\code{httpx} is declared although it is transitive.} Because the download code imports it directly, and relying on another package to pull it in is \enquote{a break waiting to happen} \ev{pyproject.toml:30-34}.

\paragraph{The Space pins majors, not exact versions.} \enquote{The Space rebuilds on its own schedule and a hard pin that stops resolving turns a working demo into a build failure months later} \ev{app/requirements.txt:6-9}. Gradio is left out entirely, because the Space builder appends its own \code{gradio[oauth,mcp]==<sdk_version>}; declaring a range as well once made the build fail with \code{ResolutionImpossible} against Gradio 6.26.0 \ev{app/requirements.txt:15-22}.

\section{Dependency history}

\begin{center}
\begin{tabularx}{\textwidth}{@{}>{\raggedright\arraybackslash}p{2.0cm}>{\raggedright\arraybackslash}p{4.2cm}L@{}}
\toprule
\textbf{Commit} & \textbf{Change} & \textbf{Recorded reason}\\
\midrule
\code{9ccd6c2} & Move to local embeddings plus Gemini (PR \#3) & The commit title calls it a \enquote{zero-cost local+Gemini stack}; \code{docs/DESIGN.md} records that the OpenAI alternative needed funded billing. The package docstring notes the swap point \enquote{was already exercised once} \ev{src/vidyarag/llm/\_\_init\_\_.py}.\\
\code{8fa9478} & LlamaIndex removed & \enquote{declared, claimed, and never used}.\\
\code{8ff561c} & \code{tenacity} removed (PR \#14) & \enquote{declared and never imported}.\\
\code{88d7969} & Single version source through hatchling (PR \#18) & Three version numbers disagreed.\\
\bottomrule
\end{tabularx}
\end{center}

\section{Where the declarations and reality disagree}
\label{sec:dep-findings}

\begin{enumerate}
  \item \textbf{Linter and formatter versions differ between local hooks and CI.} The pre-commit hooks pin ruff v0.8.4, black 24.10.0 and mypy v1.13.0, while the development group resolves to ruff 0.16.0, black 25.12.0 and mypy 1.20.2 \tagM, and CI runs the latter through \code{uv run} \ev{.pre-commit-config.yaml:35-53}, \ev{.github/workflows/ci.yml:63-70}.
  \item \textbf{The API image installs Gradio, which the API never imports.} Gradio is a runtime dependency, the Dockerfile installs the runtime set, and nothing under \code{src/vidyarag/api/} imports Gradio; only \code{app/app.py} does, and \code{app/} is not copied into the image \ev{Dockerfile:36-45}, \ev{.dockerignore}.
  \item \textbf{\code{pandas} and \code{datasets} are declared but never imported} by project code (\code{datasets} is required by RAGAS anyway). This is the same pattern that led to removing LlamaIndex and \code{tenacity}.
  \item \textbf{\code{huggingface_hub} is imported but not declared.} \code{scripts/deploy_space.py} imports \code{HfApi}; the package arrives only transitively, through Gradio and fastembed. This is exactly the situation the \code{httpx} comment warns against, in a developer tool rather than in the runtime path.
  \item \textbf{The Space and CI can run different library versions.} The lockfile governs CI and Docker; the Space resolves major-version ranges whenever it rebuilds. This is a deliberate trade-off, recorded above.
\end{enumerate}

\begin{inferred}
Item 1 means a commit that passes the local hooks can fail \code{black --check} or \code{ruff check} in CI, or the reverse, whenever the two tool versions disagree about a rule or a formatting decision. Item 2 costs image size and attack surface without adding any capability to the API.
\end{inferred}

\begin{recommended}
Keep hook versions in step with the lockfile, for example with local hooks that run \code{uv run ruff}, \code{uv run black} and \code{uv run mypy}. Move Gradio into an optional \code{demo} group and install it only for the Space. Delete \code{pandas} and \code{datasets} from the \code{eval} group, or import them. Declare \code{huggingface_hub} in a \code{deploy} group.
\end{recommended}
